\section{Theoretical Framework}
The intrinsic soft function ($S_I$) can be obtained from the QCD
factorization of a large-momentum form factor of a non-singlet light pseudo-scalar meson with constituents $\pi=\overline{q}_2\gamma_5 q_1$, with the transition current made of two quark-bilinears 
with a fixed transverse separation $\vec{b}=(\vec{n}_\perp b_\perp, 0)$,
\begin{align}\label{eq:FF}
F(b_\perp,P^z)=\langle \pi(-\vec{P})|(\overline{q}_1\Gamma q_1)(\vec{b}) (\overline{q}_2\Gamma q_2)(0)|\pi(\vec P)\rangle_c. 
\end{align}
Here $q_{1,2}$ are light quark fields of different flavors, and $\vec{P}=(\vec{0}_{\perp},P^z)$.  {To extract the soft-factor, {operators and mesonic states are chosen such that each of the four lines in Fig.~\ref{fig:formfactor} are of a different
flavor as pointed out in Ref.~\cite{Ji:2019sxk}.}. The simplest scenario would correspond to the contraction in Fig.~\ref{fig:formfactor}, which shares the same topology as the so-called connected insertion. Thus  a subscript $c$ is added on the right-hand side of Eq.~(\ref{eq:FF}). By construction, the disconnected insertion   is not relevant  in this scenario which we will adopt in this work. }


It can be shown that the form factor defined in Eq.~(\ref{eq:FF}) is factorizable into the quasi-TMDWF $\Phi$ and the intrinsic soft function $S_I$~\cite{Ji:2019sxk,Ji:2020ect}
%
\begin{align}\label{eq:FF_factorization}
&F(b_\perp,P^z)= {S_I(b_\perp)}\\ &\times {\int_0^1 dx\, dx' H(x,x',P^z)\Phi^{\dagger}(x',b_\perp,-P^z)\, \Phi(x,b_\perp,P^z)}\nonumber
\end{align}
where $H$ is {the} perturbative hard kernel.
%
The quasi-TMDWF $ \Phi$ is the Fourier transformation of the coordinate-space correlation function
\begin{align}
\label{eq:phi_def}
&\phi(z,b_\perp,P^z)=\lim_{\ell\to\infty}\frac{\phi_\ell(z,b_\perp,P^z,\ell)}{\sqrt{Z_E(2\ell,b_\perp)}},\\
&\phi_\ell(z,b_\perp,P^z,\ell)\nonumber\\
&=\Big\langle 0\Big|\overline q_1\left(\frac{z }{2}n^z +\vec b\right)\Gamma_\Phi\,{\cal W}(\vec b,\ell)q_2\left(-\frac{z}{2}n^z \right)\Big| {\pi(\vec{P})}\Big\rangle\nonumber. 
\end{align}
In the above ${\cal W}(\vec b,\ell)$ is the spacelike staple-shaped gauge link,
\begin{align}
{\cal W}(\vec b,\ell)&= {\cal P}{\rm exp} \left[ ig_s\int_{-\ell}^{z/2} \textrm{d}s\ n^z\cdot  A(n^z s+b_\perp)\right] \nonumber\\
& \times {\cal P}{\rm exp} \left[ ig_s\int_{0}^{b_\perp } \textrm{d}s\ n_\perp\cdot  A(-\ell n^z +s  n_\perp)\right] \nonumber\\
& \times {\cal P}{\rm exp} \left[ ig_s\int_{-z/2}^{-\ell} \textrm{d}s\ n^z\cdot  A(n^z s)\right],
\end{align}
$n^z$ and $n_\perp$ are the unit vectors in $z$ and  transverse directions respectively. 
$Z_E(2\ell,b_\perp)$ is the vacuum expectation value of a rectangular spacelike Wilson loop with size $2\ell \times b_\perp$ which   removes the pinch-pole singularity and Wilson-line self-energy in quasi-TMDWF~\cite{Ji:2019sxk}.

Since the UV divergence of the intrinsic soft function is multiplicative~\cite{Ji:2020ect},
the  ratio  $S_I(b_{\perp},1/a)/S_I(b_{\perp,0},1/a)$ calculable on lattice is UV
renormalization-scheme independent, where $b_{\perp,0}$ is a reference distance which is taken
small enough to be calculated perturbatively. Thus we
can obtain the result in the $\overline{\rm MS}$ scheme
through
\begin{align}\label{eq:S_ratio}
S_{I,\overline{\rm MS}}(b_\perp,\mu)= \left(\frac{S_I(b_{\perp},1/a)}{S_I(b_{\perp,0},1/a)}\right)S_{I,\overline{\rm MS}}(b_{\perp,0},\mu)
\end{align}
where $S_{I,\overline{\rm MS}}(b_{\perp,0},\mu)$ is perturbatively calculable, e.g.,
\begin{align}\label{eq:S_ratio_1loop}
S_{I,\overline{\rm MS}}(b_{{\perp}},\mu)=1-\frac{\alpha_sC_F}{\pi}\ln\frac{\mu^2 b_{{\perp}}^2}{4 e^{-2\gamma_E}}+{\cal O}(\alpha_s).
\end{align}

In the present exploratory study, we will consider only   leading
order matching in Eq. (\ref{eq:FF_factorization}), for which the perturbative
kernel is $H(x,x',P^z)=1/({2N_c})+{\cal O}(\alpha_s)$, independent of $x$ and $x'$.
Using $\phi(0,b_\perp,-P^z)=\phi(0,b_\perp,P^z)$ under parity transformation, we obtain
\begin{align}\label{eq:S_I_z=0}
S_I(b_\perp)=\frac{2N_c F(b_\perp,P^z)}{|\phi(0,b_\perp,P^z)|^2}+{\cal O}(\alpha_s, (1/P^z)^2),
\end{align}
where power corrections from finite $P^z$ are ignored. Since $P^z$ is related
to the rapidity of the meson, we henceforth replace it by the boost factor $\gamma\equiv E_{\pi}/m_{\pi}$.
Eq.~(\ref{eq:S_ratio}) can be written as
\begin{align}\label{eq:S_ratio_z=0}
S_{I,\overline{\rm MS}}(b_\perp,\mu)&=\frac{F(b_\perp,P^z)}{F(b_{\perp,0},P^z)}\frac{|\phi(0,b_{\perp,0},P^z)|^2}{|\phi(0,b_\perp,P^z)|^2}\nonumber\\
&+{\cal O}(\alpha_s, \gamma^{-2})\, .
\end{align}
%
The ratio on the right-hand side of the above expression  is independent of the renormalization scale $\mu$ since only the leading-order contribution is kept.

%To our knowledge, the quasi-TMDWF has not been studied before on the lattice. Once it is known, it
{On the other hand, the quasi-TMDWF} can be used to extract the Collins-Soper kernel $K$ using a method similar
to~\cite{Ebert:2018gzl}
\begin{align}
&K(b_\perp,\mu)=\frac{1}{\ln(P_1^z/P_2^z)}\ln\left|\frac{C(xP_2^z,\mu) \Phi_{\overline{\rm MS}}(x,b_\perp,P_1^z,\mu)}{C(xP_1^z,\mu) \Phi_{\overline{\rm MS}}(x,b_\perp,P_2^z,\mu)}\right|\label{eq:CS_kernel}\\
&\quad=\frac{1}{\ln(P_1^z/P_2^z)}\ln\left|\frac{\int_0^1 \textrm{d}x \Phi(x,b_\perp,P_1^z)}{\int_0^1 \textrm{d}x \Phi(x,b_\perp,P_2^z)}\right|+{\cal O}(\alpha_s, \gamma^{-2})\nonumber\\
&\quad=\frac{1}{\ln(P_1^z/P_2^z)}\ln\left|\frac{\phi(0,b_\perp,P_1^z)}{\phi(0,b_\perp,P_2^z)}\right|+{\cal O}(\alpha_s, \gamma^{-2})\label{eq:CS_kernel_z=0}. 
\end{align}
%
{In the second line,  again only the leading order matching kernel $C(xP^z,\mu)=1+{\cal O}(\alpha_s)$ is used. }
The  renormalization factors for $\Phi$ are cancelled.
%
The rapidity-scheme-independent CS kernel $K$ is independent of $\mu$ in this approximation because only the leading term has been kept.

While Eqs.~(\ref{eq:S_ratio}) and (\ref{eq:CS_kernel}) are exact and can be
used for precision {studies} in the future,  Eqs.~(\ref{eq:S_ratio_z=0}) and (\ref{eq:CS_kernel_z=0}) are the leading-order approximation used in this pioneering  work.


\begin{table}[htbp]
\begin{center}
\caption{\label{table:cls} Parameters used  in the numerical simulation. The first row shows the parameters of the 2+1 flavor clover fermion CLS ensemble (named A654) and the second one shows the number of the A654 configurations and valence pion mass used for this calculation.}
\begin{tabular}{ccccccc}
\hline
 $\beta$ & $L^3\times T$  &a (fm)  & $c_{sw}$  & $\kappa^{\rm sea}_l$ & $m^{\rm sea}_{\pi} $(MeV)\\
  3.34 & $24^3 \times  48$ & 0.098 &  2.06686  & 0.13675 & 333 \\
\hline
&&& $N_{cfg}$  &  $\kappa^{v}_l$  & $m^v_{\pi} $ (MeV)  \\
  &&&  864  & 0.13622   & 547 \\
 \hline
\end{tabular}
\end{center}
\end{table}
